\section{Evaluation protocol}
\label{sec:protocol}

\subsection{Splits and windows}

Temporal splits are fixed for every track: training up to 2010-12-31, validation 2011--2015, and a 2016--2023 test
period common to all tracks. Anything estimated from data (climatologies, percentile thresholds, the spell
reference cycle, and model fitting) uses the training split only.

Evaluation is rolling-origin with a 7-day stride and horizons of 1, 3, 7, 14 and 30 days. A window is scored when
its 30-day target lies entirely inside the test split, is complete and unflagged, and at least 365 days of history
precede the origin. Context is deliberately \emph{not} required to be complete: a 90\%-complete-context rule leaves
zero admissible GHCN-Bangladesh rainfall windows, which would silently delete the sparsest and most
representative part of the problem. Models therefore receive contexts containing real gaps, and each window
records its context completeness so that skill can be stratified by it. This yields 170{,}571 daily windows and
18{,}112 dekadal windows.

\subsection{Models and baselines}

The eight foundation models are run zero-shot from their published checkpoints (Section~\ref{sec:results}).
Chronos-2, Chronos-Bolt, Toto-2.0 and Moirai-2 accept missing values natively and receive the context with its gaps.
TimesFM 2.5, TiRex, TTM r2 and Sundial do not document their handling of missing values, so they receive a linearly
interpolated context, with edges filled from the nearest observation.

Each of the ten baselines is trained on its track's training split with one fixed configuration. No
hyperparameter search was run for any baseline, so every claim about baselines in this paper concerns these
configurations, not tuned ones. NHITS, DLinear and PatchTST are global models, one per track and variable, trained
for 1{,}000 steps on a 365-day input with a multi-quantile loss over the nine scored quantiles. They model
anomalies from the training climatology, with gaps zero-filled and excluded from the loss. LightGBM is one global
model per track and variable, with fixed parameters and an L2 objective on the anomaly; it stops early on the
validation split and takes its predictive quantiles from validation residuals. AutoETS, AutoTheta and AutoARIMA choose their own
order by information criterion on each window's climatological anomaly, with context gaps interpolated. The
naive references are day-of-year climatology, persistence and seasonal naive.

\subsection{Metrics}

The headline metric is the continuous ranked probability score, a strictly proper scoring rule
\citep{gneiting2007scoring}, computed here from nine predictive quantiles as twice the mean pinball loss. Point errors use each model's predictive \emph{median} and squared errors its mean.
Mixing the two is not neutral on zero-inflated rainfall. Scoring climatology on its mean rather than its median
changed its MASE at a 30-day lead from 0.58 to 0.82. It also made a quantile-loss network appear
\SI{25}{\percent} better than it is. Interval calibration is reported as coverage of the central \SI{80}{\percent} interval.

Because scaled errors differ systematically between tracks, all cross-track comparisons use skill relative to the
\emph{same track's} day-of-year climatology, estimated on the training split. A score of zero is climatology.

\subsection{Significance}

Neighbouring origins overlap heavily, so treating windows as independent samples would overstate confidence
substantially; none of the reference forecasting benchmarks reports significance tests at all. We therefore make
the forecast origin the unit of resampling, averaging stations within an origin so that cross-sectional
correlation is absorbed rather than counted as extra samples.

We use three tools. To test whether one model beats another, we apply the Diebold--Mariano test to the paired
difference in loss. Its variance is a Newey--West HAC estimate at lag $h-1$, with the Harvey--Leybourne--Newbold
small-sample correction and a $t$ reference distribution.

To put an interval around a skill score or a difference, we use a moving-block bootstrap over origins, with blocks
of 8 origins and 2{,}000 resamples. Series that are compared with each other, such as a model, its climatology
reference and a control baseline, are resampled with the same blocks. A difference-in-differences interval
therefore carries the control's uncertainty rather than treating it as fixed.

To rank models across tasks, we use the Friedman test with the Nemenyi critical difference from the studentized
range.

The primary analysis consists of 48 comparisons. For each of the 16 tasks and each of the leads 1, 7 and 30 days,
it compares the best foundation model with the best baseline, chosen as described below. The three leads are the
shortest horizon we score, one week, and the longest; leads 3 and 14 fall between them and appear in the leaderboard and figures as point estimates, with no
claim resting on them alone. Adding them would enlarge the family to 80 comparisons and make the correction
stricter for every test, without adding a horizon outside the range the three already span. All 48 comparisons are
Holm--Bonferroni corrected, and we report the corrected values.

The comparators in that family are chosen on the validation split. Every model also forecasts from origins in
2011--2015, with the same configuration or checkpoint as on the test split. Those forecasts are scored with the
same harness and training-period scales. For each task and lead, the foundation model and the baseline with the
lowest mean scaled CRPS on the validation windows are frozen. Only then is the pair tested on the 2016--2023
windows, so no test-period score enters the selection. Table~\ref{tab:comparators} names both sides.
LightGBM is the one candidate whose validation scores are not clean: it stops early on the validation split and
takes its quantiles from validation residuals, which flatters it there. We therefore repeat the selection with
LightGBM excluded.

Two further versions show how much the choice of comparators matters. The first chooses both sides on the test
windows themselves; its Diebold--Mariano $p$-values do not account for that selection. The second keeps selection
and testing apart within the test period: the pair is chosen on origins in 2016--2019 and tested on origins in
2020--2023 alone, with the same test and correction. As a reference for the loss of power from halving the
sample, the test-chosen pair is tested on the same 2020--2023 origins.

Every other statement of uncertainty in the paper (skill against climatology,
reanalysis against observations, the difference-in-differences and the spell contrasts) is a per-comparison
\SI{95}{\percent} bootstrap interval without multiplicity adjustment. We describe those as intervals that exclude
zero, not as significant results. Table~\ref{tab:evidence} lists the status of each.

\subsection{Controls}

\paragraph{Gap-matched contexts.} Bangladeshi contexts are 82--86\% observed over the year preceding an origin,
against 99--100\% for the temperate reference (Figure~\ref{fig:coverage}). To separate climate from data
availability we draw real missingness
patterns from GHCN-Bangladesh windows and impose them, right-aligned and deterministically, on temperate contexts.
Masking moves the temperate observed share to 0.808 for rainfall (Bangladesh reference 0.819) and 0.863 for mean
temperature (reference 0.863). Crucially the mask is applied to the trained baselines as well as to the foundation
models, so the control is like-for-like; masking only the model understates the remaining effect. For each
variable and lead the control is the trained baseline with the lowest test CRPS on the GHCN-Bangladesh track; for
temperature at a 30-day lead that is AutoARIMA. We summarise the effect of masking as the share of the raw median
penalty that it removes: (raw median $-$ like-for-like median) / raw median.

\paragraph{Leakage audit.} Every model's pretraining corpus was checked against its technical report rather than
its model card. No corpus names GHCN-Daily, BMD or any Bangladesh source. The station corpora that do appear,
USHCN and SubseasonalClimateUSA, are United States only, and our temperate reference is entirely European, so
neither side of the tropical-versus-temperate contrast is contaminated. Leakage-free zero-shot evaluation is not
itself novel \citep{qiao2026time}; what matters here is that it is a precondition for reading a geographic
contrast as a geographic effect.

\paragraph{Context length.} Because a sparse context is also a shorter effective context, the tropical contrast is
repeated with 2{,}048 days of history in place of 1{,}024. The control is masked identically.
